\documentclass[10pt,a4paper]{amsart}
\usepackage{amsmath,amssymb,amsthm,graphicx}
\usepackage{verbatim}
\newenvironment{reference}{\comment}{\endcomment}
%\newenvironment{reference}{}{}
\newenvironment{slogan}{\comment}{\endcomment}
\newenvironment{history}{\comment}{\endcomment}

% For commutative diagrams we use Xy-pic
\usepackage[all]{xy}

% We use 2cell for 2-commutative diagrams.
\xyoption{2cell}
% Theorem environments.
%
\theoremstyle{plain}
\newtheorem{theorem}[subsection]{Theorem}
\newtheorem{proposition}[subsection]{Proposition}
\newtheorem{lemma}[subsection]{Lemma}

\theoremstyle{definition}
\newtheorem{definition}[subsection]{Definition}
\newtheorem{example}[subsection]{Example}
\newtheorem{exercise}[subsection]{Exercise}
\newtheorem{situation}[subsection]{Situation}

\theoremstyle{remark}
\newtheorem{remark}[subsection]{Remark}
\newtheorem{remarks}[subsection]{Remarks}

\numberwithin{equation}{subsection}

% Macros
%
\def\lim{\mathop{\mathrm{lim}}\nolimits}
\def\colim{\mathop{\mathrm{colim}}\nolimits}
\def\Spec{\mathop{\mathrm{Spec}}}
\def\Hom{\mathop{\mathrm{Hom}}\nolimits}
\def\Ext{\mathop{\mathrm{Ext}}\nolimits}
\def\SheafHom{\mathop{\mathcal{H}\!\mathit{om}}\nolimits}
\def\SheafExt{\mathop{\mathcal{E}\!\mathit{xt}}\nolimits}
\def\Sch{\mathit{Sch}}
\def\Mor{\mathop{\mathrm{Mor}}\nolimits}
\def\Ob{\mathop{\mathrm{Ob}}\nolimits}
\def\Sh{\mathop{\mathit{Sh}}\nolimits}
\def\NL{\mathop{N\!L}\nolimits}
\def\CH{\mathop{\mathrm{CH}}\nolimits}
\def\proetale{{pro\text{-}\acute{e}tale}}
\def\etale{{\acute{e}tale}}
\def\QCoh{\mathit{QCoh}}
\def\Ker{\mathop{\mathrm{Ker}}}
\def\Im{\mathop{\mathrm{Im}}}
\def\Coker{\mathop{\mathrm{Coker}}}
\def\Coim{\mathop{\mathrm{Coim}}}

% Boxtimes
%
\DeclareMathSymbol{\boxtimes}{\mathbin}{AMSa}{"02}

%
% Macros for moduli stacks/spaces
%
\def\QCohstack{\mathcal{QC}\!\mathit{oh}}
\def\Cohstack{\mathcal{C}\!\mathit{oh}}
\def\Spacesstack{\mathcal{S}\!\mathit{paces}}
\def\Quotfunctor{\mathrm{Quot}}
\def\Hilbfunctor{\mathrm{Hilb}}
\def\Curvesstack{\mathcal{C}\!\mathit{urves}}
\def\Polarizedstack{\mathcal{P}\!\mathit{olarized}}
\def\Complexesstack{\mathcal{C}\!\mathit{omplexes}}
% \Pic is the operator that assigns to X its picard group, usage \Pic(X)
% \Picardstack_{X/B} denotes the Picard stack of X over B
% \Picardfunctor_{X/B} denotes the Picard functor of X over B
\def\Pic{\mathop{\mathrm{Pic}}\nolimits}
\def\Picardstack{\mathcal{P}\!\mathit{ic}}
\def\Picardfunctor{\mathrm{Pic}}
\def\Deformationcategory{\mathcal{D}\!\mathit{ef}}
\usepackage{cleveref}
\crefname{originallemma}{original lemma}{original lemmas}
\Crefname{originallemma}{Original lemma}{Original lemmas}
\makeatletter
\newlabel{categories-lemma-left-dual}{{0FFQ}{source}{}{sourcepointer-1}{}}
\newlabel{categories-lemma-left-dual@cref}{{[originallemma][0][]0FFQ}{[1][0][]source}{}{}{}}
\makeatother
\crefname{originallemma}{original lemma}{original lemmas}
\Crefname{originallemma}{Original lemma}{Original lemmas}
\makeatletter
\newlabel{categories-lemma-yoneda}{{001P}{source}{}{sourcepointer-2}{}}
\newlabel{categories-lemma-yoneda@cref}{{[originallemma][0][]001P}{[1][0][]source}{}{}{}}
\makeatother
\crefname{originalreference}{original reference}{original references}
\Crefname{originalreference}{Original reference}{Original references}
\makeatletter
\newlabel{lemma-chow-groups-representable}{{0FGM}{source}{}{sourcepointer-3}{}}
\newlabel{lemma-chow-groups-representable@cref}{{[originalreference][0][]0FGM}{[1][0][]source}{}{}{}}
\makeatother
\crefname{originalreference}{original reference}{original references}
\Crefname{originalreference}{Original reference}{Original references}
\makeatletter
\newlabel{lemma-dual-general}{{0FGJ}{source}{}{sourcepointer-4}{}}
\newlabel{lemma-dual-general@cref}{{[originalreference][0][]0FGJ}{[1][0][]source}{}{}{}}
\makeatother
\title{Manin's criterion for an isomorphism of motives}
\author{The Stacks Project Authors}
\date{stacks-project:weil.tex:local-snapshot-20261012; GNU FDL 1.2 or later}
\begin{document}
\maketitle
\noindent\textit{Editorial note.} All delivered mathematical text is the author's own, copied from the pinned Stacks Project chapter file weil.tex; no statement or proof step was rewritten, reordered or removed.\par\medskip
\noindent\textit{Editorial note.} The counted proof is the author's complete argument: the reduction to arbitrary objects through the Karoubian envelope and the twist, the translation into morphism groups out of the unit, and the final use of left duals and the Yoneda lemma.\par\medskip
\noindent\textit{Editorial note.} The label lemma-manin names the counted statement, and the proof environment immediately following it in the source is delivered verbatim as its proof.\par\medskip
\noindent\textit{Editorial note.} The counted statement is carried by the author's own lemma environment in the chapter source.\par\medskip
\noindent\textit{Editorial note.} Pre-open gate: python operations/agent-authoring/tools/precard\_gate.py --source weil.tex --label lemma-manin --start 1326 --end 1356 --local-id 28149 returned RESULT=OK (L28149 free, occupancy FREE, public interval clear). Pre-assembly tag self-check: canonical tag resolved and not present in reference/public-items.jsonl.\par\medskip
\noindent\textit{Editorial note.} The statement and proof are copied from the pinned chapter file; the cited inputs are kept as external pointers with their canonical Stacks tags.\par\medskip
\noindent\textit{Allowed standard prerequisites (editorial).}\begin{itemize}
\item Objects of the category of motives are summands of twisted motives of schemes. The construction of the category of motives as the Karoubian envelope of the category of correspondences, so that every object is a summand of a twisted motive of a smooth projective scheme. Reduces the assumption from schemes to arbitrary objects, using that taking Chow groups commutes with the twist up to a shift.
\item Chow groups as morphism groups from the unit. Same chapter, the lemma representing Chow groups as morphism groups out of the unit. Converts the isomorphism on Chow groups into an isomorphism of morphism groups out of the unit after tensoring with any object.
\item Existence of left duals and the Yoneda lemma. Same chapter, the lemma that every object of the category of motives has a left dual, and the categorical lemma that a morphism inducing isomorphisms on all morphism groups is an isomorphism. Upgrades the isomorphism with the unit on the left to an isomorphism of all morphism groups and concludes by the Yoneda lemma.
\end{itemize}\medskip
\ifdefined\hypertarget\hypertarget{sourcepointer-1}{}\fi
\noindent\textit{Source reference.} Original lemma 0FFQ: sources/stacks-chow/tags.txt, source lines 19295--19295. Prerequisite pointer invoked at proof line(s) 1353; it names a standard fact of a sibling chapter that the counted argument uses, and every step of the counted proof built on it is retained verbatim in the delivered body.\par\medskip
\ifdefined\hypertarget\hypertarget{sourcepointer-2}{}\fi
\noindent\textit{Source reference.} Original lemma 001P: sources/stacks-chow/tags.txt, source lines 70--70. Prerequisite pointer invoked at proof line(s) 1355; it names a standard fact of a sibling chapter that the counted argument uses, and every step of the counted proof built on it is retained verbatim in the delivered body.\par\medskip
\ifdefined\hypertarget\hypertarget{sourcepointer-3}{}\fi
\noindent\textit{Source reference.} Original reference 0FGM: sources/stacks-weil/pinned-original-source.tex, source lines 1306--1306. Prerequisite pointer invoked at proof line(s) 1341; it names a standard fact of the same chapter that the counted argument uses, and every step of the counted proof built on it is retained verbatim in the delivered body.\par\medskip
\ifdefined\hypertarget\hypertarget{sourcepointer-4}{}\fi
\noindent\textit{Source reference.} Original reference 0FGJ: sources/stacks-weil/pinned-original-source.tex, source lines 1238--1238. Prerequisite pointer invoked at proof line(s) 1348; it names a standard fact of the same chapter that the counted argument uses, and every step of the counted proof built on it is retained verbatim in the delivered body.\par\medskip
\begin{lemma}[Manin]
\label{lemma-manin}
Let $k$ be a base field. Let $c : M \to N$ be a morphism of motives.
If for every smooth projective scheme $X$ over $k$ the map
$c \otimes 1 : M \otimes h(X) \to N \otimes h(X)$ induces an isomorphism on
Chow groups, then $c$ is an isomorphism.
\end{lemma}

\begin{proof}
Any object $L$ of $M_k$ is a summand of $h(X)(m)$ for some smooth projective
scheme $X$ over $k$ and some $m \in \mathbf{Z}$. Observe that the Chow groups
of $M \otimes h(X)(m)$ are the same as the Chow groups of $M \otimes h(X)$
up to a shift in degrees. Hence our assumption implies
that $c \otimes 1 : M \otimes L \to N \otimes L$ induces an isomorphism on
Chow groups for every object $L$ of $M_k$. By
Lemma \ref{lemma-chow-groups-representable}
we see that
$$
\Hom_{M_k}(\mathbf{1}, M \otimes L) \to
\Hom_{M_k}(\mathbf{1}, N \otimes L)
$$
is an isomorphism for every $L$. Since every object of $M_k$ has a left dual
(Lemma \ref{lemma-dual-general}) we conclude that
$$
\Hom_{M_k}(K, M) \to \Hom_{M_k}(K, N)
$$
is an isomorphism for every object $K$ of $M_k$, see
Categories, Lemma \ref{categories-lemma-left-dual}.
We conclude by the Yoneda lemma
(Categories, Lemma \ref{categories-lemma-yoneda}).
\end{proof}

\end{document}
